\documentclass[10pt]{article}
\usepackage[margin=20mm]{geometry}
\usepackage{amsmath,amssymb,mathtools,fontspec,unicode-math,xurl,hyperref}
\IfFontExistsTF{DejaVu Serif}{\setmainfont{DejaVu Serif}}{\setmainfont{Latin Modern Roman}}
\setmathfont{latinmodern-math.otf}
\setlength{\parindent}{0pt}\setlength{\parskip}{6pt}\sloppy
\title{Technical verification of
degree theory for expanding solitons}
\author{}
\date{}
\begin{document}
\maketitle
Abstract: This report examines two specific steps in Bamler and Chen’s arXiv manuscript 2305.03154v3: the drift estimate for generalized-cone deformations and the comparison of finite-dimensional analytic slices used to define the degree. Explicit geometric and finite-dimensional calculations identify their precise scope. The reviewed object is the author manuscript, not the later journal version. Scientific adjudication and independent acceptance remain pending.\par
\section*{Identity, source basis, and scope}
The reviewed paper is Richard H. Bamler and Eric Chen, Degree theory for 4-dimensional asymptotically conical gradient expanding solitons, arXiv:2305.03154v3, submitted 23 January 2025, with internal date 27 January 2025 [BC]. The frozen PDF has 122 physical pages; throughout this report page numbers refer to those one-based physical pages, which agree with the manuscript’s printed pagination. The associated journal DOI is 10.1002/cpa.70024 (https://doi.org/10.1002/cpa.70024). The preserved publisher metadata describes a different, 148-page journal article; byte or mathematical equivalence has not been established.\par
The inspected scope consists of Definitions 3.1, 3.5, 3.8, and 3.30; Lemma 3.31 and its ordinary-cone applicability; the explicitly gradient application in Lemma 8.11; and the slice-comparison step in Lemma 10.39 with its relevant admissibility conditions and dependency on Theorem 10.49. A saved hyperbolic test of Proposition 5.33 was checked against the standing asymptotically conical convention and is not promoted to a defect. The full compactness, elliptic, gauge, index-change, and external-reference arguments have not been independently re-established. Detailed inspected and uninspected items are recorded in `coverage.json`.\par
\section*{Generalized cones and the drift estimate}
Equation reference: eq:gencone. Fix a finite integer \(k^*\geq 5\); one may take \(k^*=30\). Definition 3.1 allows \[
 \gamma=dr^2+r(dr\otimes\beta+\beta\otimes dr)+r^2h,\] where \(h\) and \(\beta\) are smooth tensors on the compact link and the displayed tensor is positive definite. The tangent space consists of homogeneous symmetric tensors with zero \(dr^2\) component. An ordinary cone has \(\beta=0\). Definition 3.8 uses a specified soliton vector field \(V\), and requires \(\iota^*V=-r\partial_r/2\) and conical asymptotics. It does not require \(V\) to be gradient. These distinctions are explicit on pp. 20–23.\par
Equation reference: eq:claimed. Problem 1 (The second estimate of Lemma 3.31 fails for generalized cones).  For an arbitrary representative \((g,V,\gamma)\) in the full space \(\mathcal M^{k^*}\), Lemma 3.31, p. 29, asserts \[
 \left|\nabla^m\nabla_V T(\gamma')\right|\circ\iota
 \leq C r^{-2-m}\|\gamma'\|,\qquad 0\leq m\leq k^*-5,\] for every \(\gamma'\in T\operatorname{GenCone}^{k^*}(N)\). Already its \(m=0\) case is false.\par
\subsection*{An admissible complete soliton}
Let \(M=\mathbb R^4\), \(N=S^3\), and \(g\) be the Euclidean metric. Let \[J=\begin{pmatrix}0&-1&0&0\\1&0&0&0\\0&0&0&-1\\0&0&1&0\end{pmatrix},
 \qquad V(x)=\tfrac12(-I+J)x .\] Thus \(J^\top=-J\), \(J^2=-I\), and \(J^\top J=I\). Both \(g\) and \(V\) are smooth and complete: the vector-field flow is \(x(t)=e^{-t/2}e^{tJ/2}x(0)\) for all real \(t\). Since \(\operatorname{Ric}_g=0\) and \(\mathcal L_Vg=-g\), the soliton equation \(\operatorname{Ric}_g+\tfrac12\mathcal L_Vg+\tfrac12g=0\) holds.\par
Equation reference: eq:embedding. Define the end embedding \[
 \iota(r,\theta)=\frac{r}{\sqrt2}e^{-(\log r)J}\theta,\qquad r\geq1,\quad\theta\in S^3.\] It is a smooth embedding onto \(\{|x|\geq1/\sqrt2\}\), and the complement of each truncated end is compact. The link \(S^3\) admits positive scalar curvature, and the identity cover of \(\mathbb R^4\) has \(H_2=H_1=0\). Consequently Definition 3.5’s three ensemble conditions hold.\par
Equation reference: eq:pullback. For the round unit-sphere metric \(h_0\) and \(\eta_\theta(w)=\langle J\theta,w\rangle\), direct differentiation gives \[
 \gamma:=\iota^*g
 =dr^2-\frac r2(dr\otimes\eta+\eta\otimes dr)+\frac{r^2}{2}h_0.\] Here \(|\eta|_{h_0}=1\), so for \(\beta=-\eta/2\) and \(h=h_0/2\), \(|\beta|_h^2=1/2<1\). Thus [eq:pullback] (\#eq:pullback) is positive definite, \(\gamma(\partial_r,\partial_r)=1\), and \(\mathcal L_{r\partial_r/2}\gamma=\gamma\); it is an admissible generalized cone. It is not an ordinary cone in this fixed radial parametrization. Moreover, \[d\iota(r\partial_r)=(I-J)x,\qquad
 d\iota(-r\partial_r/2)=V(x).\] All three conditions of Definition 3.8 hold: the soliton equation is exact, the vector-field condition is exact, and \(\iota^*g-\gamma=0\), so every required asymptotic derivative is zero.\par
\subsection*{The contradictory drift calculation}
Equation reference: eq:drift. Take the smooth tangent direction \(\gamma'=\gamma-dr^2\in T\operatorname{GenCone}^{k^*}(S^3)\). Its link norm is finite and independent of \(r\). The cutoff in Definition 3.30 equals one for \(r\geq3\). Because \(r=\sqrt2\,|x|\), pushing forward \(dr^2\) gives \(2xx^\top/|x|^2\). On that end, therefore, \[H:=T(\gamma')=I-2ee^\top,\qquad e=x/|x|.\] The Euclidean covariant derivative is ordinary componentwise differentiation. The radial part of \(V\) annihilates the degree-zero tensor \(H\), whereas its rotational part gives \[
 \nabla_VH=-\bigl((Je)e^\top+e(Je)^\top\bigr),\qquad
 |\nabla_VH|_g^2=2.\] Indeed \(e\) and \(Je\) are orthogonal unit vectors, so the two rank-one terms have squared Frobenius norms one and zero cross product. Equation [eq:claimed] (\#eq:claimed) would imply \(\sqrt2\leq C r^{-2}\|\gamma'\|\) for every sufficiently large \(r\), which is impossible. No approximation or empirical inference enters this contradiction.\par
The saved SymPy 1.13.3 computation verifies the matrix identities, soliton residual, pullback radial normalization, and squared norm in [eq:drift] (\#eq:drift). It is supporting algebra, not a replacement for the completeness, topology, embedding, and asymptotic checks above.\par
\subsection*{The failed step, counterevidence, and a supported restriction}
The proof on p. 29 begins with \(\nabla^\gamma_{\partial_r}\gamma'=0\). Homogeneity supplies a Lie-derivative identity; it does not supply this covariant-derivative identity for a general metric [eq:gencone] (\#eq:gencone). In the example, \(\iota^*g=\gamma\) exactly, so no comparison-connection error can cancel the nonzero leading term. This directly identifies the error in the passage to equation (3.32).\par
Changing \(V\) to the Gaussian gradient field is not an admissible way to dismiss the example: \(V\) and the end embedding are part of the original object. Remark 3.9, p. 23, explicitly distinguishes a specified gradient vector field from a metric that admits some other gradient soliton vector field. Conversely, this example does not contradict a version of the lemma restricted to ordinary asymptotic cones.\par
A supported corrected statement is obtained by adding \(\gamma\in\operatorname{Cone}^{k^*}(N)\) to Lemma 3.31, with the other finite-regularity hypotheses unchanged. On an ordinary cone, a homogeneous tangent tensor \(\gamma'\) has \(\nabla^\gamma_{\partial_r}\gamma'=0\). The asymptotic estimate (3.34) then gives a connection difference of order \(r^{-3}\) and, after contraction with \(r\partial_r\), order \(r^{-2}\). Differentiating through the stated ranges gives precisely [eq:claimed] (\#eq:claimed); the first estimate follows from homogeneity and the same comparison. The compact cutoff region is absorbed into \(C\). Thus the original proof works under this additional hypothesis.\par
This is assessed as a Minor scope defect. The explicitly inspected application in Lemma 8.11, pp. 74–76, fixes a gradient base soliton, whose cone is ordinary by Remark 3.9, so that application is covered by the corrected statement. No failure of a central degree theorem is established by this counterexample, and the complete dependency network has not been re-proved.\par
\section*{The common analytic slice in Lemma 10.39}
Problem 2 (An unsupported dimension bound for a common slice).  The last paragraph of the proof of Lemma 10.39, p. 108, takes two arbitrary finite-dimensional real-analytic submanifolds \(Q_1,Q_2\) through a cone \(\gamma\), each satisfying the transversality condition (10.4). It asserts that smaller neighborhoods in both are contained in a real-analytic submanifold \(Q_3\) of dimension \(\dim Q_1+\dim Q_2\). This dimension assertion is false, including for admissible slices in the paper’s setting.\par
\subsection*{Construction and full transversality condition}
Equation reference: eq:slice. Use the standard ensemble \((\mathbb R^4,S^3,\iota)\), with \(\iota(r,\theta)=r\theta\), \(k^*=30\), and \(\alpha=1/2\). Its topological conditions hold as above. Let \(\gamma_0=dr^2+r^2h_0\) be the Euclidean cone. Its scalar curvature is zero, so \(\gamma_0\in\operatorname{Cone}_{R\geq0}^{k^*}(S^3)\). Choose five nonzero smooth functions \(\phi_j\) on \(S^3\) with pairwise disjoint supports and \(0\leq\phi_j\leq1\), and set \[e_j=r^2\phi_j h_0,\qquad 1\leq j\leq5 .\] They are linearly independent elements of \(T\operatorname{Cone}^{k^*}\). For sufficiently small \(x,y\), define \[\begin{aligned}
 F_1(x,y)&=\gamma_0+xe_1+ye_2,\\
 F_2(x,y)&=\gamma_0+xe_1+ye_2+x^2e_3+xye_4+y^2e_5. \end{aligned}\] Both maps are polynomial into the \(C^{30}\) Banach space and hence real-analytic. Their projections onto the first two coefficients are the identity, so they are embeddings, with common tangent space \(E=\operatorname{span}(e_1,e_2)\) at \(\gamma_0\). They are orientable by \((x,y)\). Disjoint supports and \(|x|,|y|<\varepsilon<1/2\) make the angular metric positive definite. Their images are therefore contained in \(\operatorname{GenCone}\) (in fact in \(\operatorname{Cone}\)). The definitions before Lemma 10.36, pp. 105–106, do not require all of \(Q_i\) to have nonnegative scalar curvature. The required intersection with \(\operatorname{Cone}_{R\geq0}\) contains \(\gamma_0\).\par
To verify (10.4), it is not enough to assert transversality at this one cone. Proposition 4.2(a), p. 33, gives the singleton Gaussian fiber over \(\gamma_0\) in this ensemble. The Gaussian operator has zero kernel, as also used on p. 110: each Cartesian component of a decaying kernel tensor solves \((\Delta+x\cdot\nabla/2)u=0\), and the maximum principle on expanding balls implies \(u=0\). By Proposition 10.1(c,d), p. 97, \(D\Pi\) is Fredholm of index zero and has zero kernel at this fiber, hence is surjective. Thus \(E\) satisfies (10.32) at every nonnegative-scalar gradient point over \(\gamma_0\).\par
The second assertion of Lemma 10.31, pp. 103–104, now permits each embedded surface to be shrunk about \(\gamma_0\) so that (10.4) holds at every relevant gradient fiber over the entire smaller surface. Its properness argument excludes additional nearby fibers escaping the transverse neighborhood. Shrinking changes none of the derivatives below. This admissibility check uses the stated Gaussian rigidity and properness results as internal dependencies; it does not independently re-prove those global theorems.\par
\subsection*{Why dimension four is impossible}
Suppose a \(C^2\) embedded submanifold \(P\) of the ambient Banach space contains neighborhoods of \(\gamma_0\) in both surfaces. Put \(T=T_{\gamma_0}P\). Its second fundamental form may be regarded as a symmetric bilinear map \[B:T\times T\longrightarrow \mathcal B/T,\] where \(\mathcal B=T_{\gamma_0}\operatorname{GenCone}\); no ambient Hilbert structure is needed. The first parameterization is affine, so its second derivatives vanish. Consequently \(B\) vanishes on \(E\times E\). The two parameterizations have the same first derivative at the origin. The second derivatives of [eq:slice] (\#eq:slice) must therefore satisfy \[[2e_3]=[e_4]=[2e_5]=0\quad\hbox{in }\mathcal B/T.\] Thus \(T\) contains \(e_1,\ldots,e_5\), giving \(\dim P\geq5\). In particular no real-analytic common containing manifold of dimension \(2+2=4\) exists, however small the neighborhoods are.\par
The preserved finite calculation records the common tangent matrix, the three second derivatives, and the matrix \(\operatorname{diag}(1,1,2,1,2)\) of their spanning vectors. Its determinant is \(4\) and rank is \(5\). The argument using \(B\) explains why that rank obstructs a containing manifold; a rank computation alone would not do so.\par
\subsection*{What the counterexample changes}
Lemma 10.39 establishes independence of the degree from the chosen slice. Theorem 10.49, p. 109, explicitly invokes it. This places the disputed sentence on a central proof path. Nevertheless, the nested-slice comparison (10.40) uses only a suitable finite-dimensional common slice; it does not use the asserted dimension equality. For the example above, a small open subset of \(\gamma_0+\operatorname{span}(e_1,\ldots,e_5)\) is such a common slice. Its tangent contains \(E\), and Lemma 10.31 makes it satisfy (10.4) after shrinking. Both nested comparisons can then be made.\par
Accordingly the demonstrated defect is assessed as Minor: the local dimension bound is false, but the bound itself has not been shown to be essential to the degree comparison. Neither different values of the two degrees nor failure of Lemma 10.39 or Theorem 10.49 has been established. A general justification of a common analytic refinement for arbitrary curved slices is not supplied by this bounded check; that separate question remains unresolved, rather than being counted as a demonstrated defect.\par
There is also a supported restricted repair. If \(Q_i\) are open subsets of affine slices \(\gamma+E_i\), use a neighborhood in \(\gamma+(E_1+E_2)\). It has dimension at most \(\dim E_1+\dim E_2\), contains sufficiently small parts of both slices, and preserves transversality at \(\gamma\) because its tangent contains \(E_1\). Lemma 10.31 supplies (10.4) on a smaller neighborhood. This proves the affine-slice case only. It is not presented as a universal repair for arbitrary analytic submanifolds.\par
\section*{Counterevidence and positive checks}
\subsection*{The hyperbolic candidate is outside the standing scope}
A preserved calculation used the complete hyperbolic metric \(g=6z^{-2}I\) on \(\{z>0\}\subset\mathbb R^4\), with \(f=-1\), \(V=0\), \(\operatorname{Ric}=-g/2\), and \(S:=\sup_M f=-1\) (Definition 5.5). In its formal weight, \((S-f+1)^{-a/2}=1\), the constant section \(u=1\) is a nonzero solution of \(\Delta u-\nabla_{\nabla f}u=0\). These algebraic statements are correct. They do not contradict Proposition 5.33 in its intended setting. The global convention on p. 22 restricts the paper to asymptotically conical solitons, and p. 43 explicitly restricts the weighted norms to that setting, with \(f\) tending quadratically to minus infinity. Hyperbolic space is not asymptotically conical and its constant potential does not meet that convention.\par
The original saved result’s standalone-scope interpretation is therefore withdrawn in this continuation. The original output is retained unaltered as historical evidence; it is not counted as a present finding. No weakening of the theorem’s actual standing assumptions is used to manufacture a counterexample.\par
\subsection*{Gaussian signs and the weighted operator}
For the Euclidean Gaussian expander, put \(q=|x|^2\) and \(L=\Delta+(x/2)\cdot\nabla\). Direct differentiation gives, for smooth \(u\), \[-e^{q/8}L(e^{-q/8}u)
   =-\Delta u+\left(\frac{q}{16}+\frac n4\right)u .\] This is the appropriate conjugation from the weighted measure \(e^{q/4}\,dx\). The oscillator \(-\Delta+q/16\) has lower bound \(n/4\): on compactly supported smooth functions this follows by expanding \(\sum_j\|(\partial_j+x_j/4)u\|_2^2\geq0\). The Gaussian attains that lower bound after closure. Thus \(-L\) has bottom eigenvalue \(n/2=2\) for \(n=4\) in this weighted realization, and neither negative nor zero eigenvalues. The saved symbolic conjugation residual is exactly zero. This checks the sign and the Gaussian base case; it does not verify the general index-change theorem.\par
\section*{Assessment and remaining limits}
The two findings concern different local assertions. Problem 1 disproves an estimate in the nongradient generalized-cone space and gives an ordinary-cone restriction supported by the existing proof. Problem 2 disproves a dimension bound while preserving the comparison for the explicit example and for affine slices. Their combination does not establish a failure of the main expander degree.\par
Item | Issue judgment | Repair verification | Adoption / remaining work\par
Problem 1 | Supported; Minor | Ordinary-cone restriction justified | Not requested; other uses not exhaustively traced\par
Problem 2 | Supported; Minor | Explicit five-dimensional and affine cases only | Not requested; arbitrary curved slices unresolved\par
Hyperbolic test | Withdrawn candidate | Standing-scope clarification | No correction proposed\par
Gaussian check | No issue in this check | Exact sign calculation | No general index certification\par
The preserved arXiv metadata retrieved on 29 October 2026 lists v3 as the latest public arXiv version. The bounded retrieved metadata contains no correction-link candidate; the publisher landing request returned HTTP 403. This is not proof that no erratum exists. The 148-page journal version was not acquired, and neither survival nor correction of these author-manuscript issues in that version has been verified. No author code or empirical experiment is needed for the two explicit analytic counterexamples. No author adoption or human adjudication has been recorded.\par
\section*{Supporting checks}
The reviewed source is `papers/0437/0437.pdf`, 1,496,247 bytes, SHA256\par
> `6eb39594c8d12363308d88c203206c83628a99ce3650c73b06060c4c120621c4`.\par
Its unchanged manifest is 26,229 bytes, SHA256\par
> `4e4cf45b145603dc7cb3c35f124f1044bcf1541cb81db992b9498ab7ea30e48a`.\par
All paths in this appendix are relative to the assigned staging root. The machine-readable provenance supplies current byte/hash descriptors for the report, coverage ledger, findings, source views, calculations, compiler/dependency payloads, logs, and final report-page inspections.\par
The interrupted original response said that draft work existed, but no usable manuscript text was recovered from the assigned files or referenced saved output locations. This does not establish absence from every historical conversation. The present exposition is new continuation work. The script is `independent\_calculations.py`; its result is `calculations/independent\_results.json`. The corresponding finite-slice script is `finite\_slice\_calculations.py`; its receipt records success at 04:55:42–04:55:43 UTC and its result is `slice\_calculations/results.json`. Those executions belong to the original initial-review session and were not rerun for packaging. Historical candidate labels P01 and P03 correspond to Problems 1 and 2; P02 is the withdrawn hyperbolic interpretation.\par
The present initial-review continuation completed the analytical admissibility checks, counterevidence, bounded corrections, and report. These roles are not independent acceptance.\par
Scientific decision and adjudication remain pending, and human reviews are empty. The final exact-source offline compiler run, complete before/after dependency payloads, PDF parsing, and genuine per-page viewing facts are bound in the delivered provenance. Rendering alone is never recorded as viewing.\par
9 R. H. Bamler and E. Chen, Degree theory for 4-dimensional asymptotically conical gradient expanding solitons, arXiv:2305.03154v3, 23 January 2025, internal date 27 January 2025, 122 pp., https://arxiv.org/abs/2305.03154v3. Associated journal DOI: https://doi.org/10.1002/cpa.70024 (journal text not reviewed).\par
\end{document}
